\section{Conclusion}
\label{sec:conclusion}

We present \glsf{peach}, an in-context framework that conditions a graph network simulator on a handful of point cloud sequences to predict mesh-based dynamics for unseen materials.
% Relying only on point cloud contexts further allows \gls{peach} to simulate real-world scenes without requiring either explicit parameter estimation or any test-time optimization.
By conditioning on point clouds instead of e.g., meshes, context data can be obtained easily from real world scenes.
This allows \gls{peach} to simulate real-world scenes without explicit parameter estimation or test-time optimization, and we demonstrate zero-shot sim-to-real transfer on a challenging trampoline scene.
% As such, it facilitates efficient test-time real-to-sim transfer, in turn enabling downstream tasks such as planning and control on deformable objects directly from raw sensor data.
Across four simulated datasets, \gls{peach} outperforms all learned point cloud encoders and performs on par with the mesh-based \gls{mango}, while being about $140$ times faster than test-time optimization.
% Our method combines a spatio-temporal encoder that treats context sequences as point clouds in 4D space-time with a trajectory-based learned simulator conditioned on the resulting encoded material parameters.
Crucial to our method are our novel spatio-temporal encoder that treats context sequences as point clouds in 4D space-time, and auxiliary supervision in the form of direct parameter regression and prediction of \gls{sdf} values.

\textbf{Limitations and future work.} 
\gls{peach} assumes access to an initial scene geometry and requires context trajectories that share physical properties with the target process. 
Generating an initial mesh from observation would relax the geometry assumption, while accumulating context online during, e.g., robot execution would handle material drift across a process.
Additionally, \gls{peach} currently predicts the full trajectory in a single forward pass, which may impose memory constraints for long sequences and high-resolution meshes. 
Predicting sparse keyframes and training a separate network to interpolate between them is a promising way to decouple memory cost from horizon length.
Further extensions could include other sensor modalities such as temperature readings and touch sensors for more accurate parameter estimation.
We discuss broader impact in Appendix~\ref{app:broader_impact}.



% We introduced \glsf{peach}, an in-context learning framework that adapts graph network simulators to unseen material properties using only sequences of point cloud observations.
% By decoupling material inference from mesh-based dynamics prediction, \gls{peach} extends prior mesh-dependent approaches to more realistic observation settings while retaining the efficiency and accuracy of trajectory-level graph network simulators.
% Our experiments on simulated deformation environments demonstrate that \gls{peach} achieves performance on par with mesh-based context methods, despite relying on significantly sparser input data.

% While this work takes an important step toward real-world adaptive simulation, all evaluations are conducted using simulated point clouds. Thus, real-world sensing effects such as noise, occlusions, and domain shift are not yet considered.
% Furthermore, inference still relies on a canonical simulation mesh, and we do not explicitly study geometric mismatch between simulated and observed systems.

